% Options for packages loaded elsewhere
% Options for packages loaded elsewhere
\PassOptionsToPackage{unicode}{hyperref}
\PassOptionsToPackage{hyphens}{url}
\PassOptionsToPackage{dvipsnames,svgnames,x11names}{xcolor}
%
\documentclass[
  12pt,
  letterpaper,
  DIV=11,
  numbers=noendperiod]{scrartcl}
\usepackage{xcolor}
\usepackage{amsmath,amssymb}
\setcounter{secnumdepth}{5}
\usepackage{iftex}
\ifPDFTeX
  \usepackage[T1]{fontenc}
  \usepackage[utf8]{inputenc}
  \usepackage{textcomp} % provide euro and other symbols
\else % if luatex or xetex
  \usepackage{unicode-math} % this also loads fontspec
  \defaultfontfeatures{Scale=MatchLowercase}
  \defaultfontfeatures[\rmfamily]{Ligatures=TeX,Scale=1}
\fi
\usepackage{lmodern}
\ifPDFTeX\else
  % xetex/luatex font selection
  \setmainfont[]{Times New Roman}
\fi
% Use upquote if available, for straight quotes in verbatim environments
\IfFileExists{upquote.sty}{\usepackage{upquote}}{}
\IfFileExists{microtype.sty}{% use microtype if available
  \usepackage[]{microtype}
  \UseMicrotypeSet[protrusion]{basicmath} % disable protrusion for tt fonts
}{}
\makeatletter
\@ifundefined{KOMAClassName}{% if non-KOMA class
  \IfFileExists{parskip.sty}{%
    \usepackage{parskip}
  }{% else
    \setlength{\parindent}{0pt}
    \setlength{\parskip}{6pt plus 2pt minus 1pt}}
}{% if KOMA class
  \KOMAoptions{parskip=half}}
\makeatother
% Make \paragraph and \subparagraph free-standing
\makeatletter
\ifx\paragraph\undefined\else
  \let\oldparagraph\paragraph
  \renewcommand{\paragraph}{
    \@ifstar
      \xxxParagraphStar
      \xxxParagraphNoStar
  }
  \newcommand{\xxxParagraphStar}[1]{\oldparagraph*{#1}\mbox{}}
  \newcommand{\xxxParagraphNoStar}[1]{\oldparagraph{#1}\mbox{}}
\fi
\ifx\subparagraph\undefined\else
  \let\oldsubparagraph\subparagraph
  \renewcommand{\subparagraph}{
    \@ifstar
      \xxxSubParagraphStar
      \xxxSubParagraphNoStar
  }
  \newcommand{\xxxSubParagraphStar}[1]{\oldsubparagraph*{#1}\mbox{}}
  \newcommand{\xxxSubParagraphNoStar}[1]{\oldsubparagraph{#1}\mbox{}}
\fi
\makeatother

\usepackage{color}
\usepackage{fancyvrb}
\newcommand{\VerbBar}{|}
\newcommand{\VERB}{\Verb[commandchars=\\\{\}]}
\DefineVerbatimEnvironment{Highlighting}{Verbatim}{commandchars=\\\{\}}
% Add ',fontsize=\small' for more characters per line
\usepackage{framed}
\definecolor{shadecolor}{RGB}{241,243,245}
\newenvironment{Shaded}{\begin{snugshade}}{\end{snugshade}}
\newcommand{\AlertTok}[1]{\textcolor[rgb]{0.68,0.00,0.00}{#1}}
\newcommand{\AnnotationTok}[1]{\textcolor[rgb]{0.37,0.37,0.37}{#1}}
\newcommand{\AttributeTok}[1]{\textcolor[rgb]{0.40,0.45,0.13}{#1}}
\newcommand{\BaseNTok}[1]{\textcolor[rgb]{0.68,0.00,0.00}{#1}}
\newcommand{\BuiltInTok}[1]{\textcolor[rgb]{0.00,0.23,0.31}{#1}}
\newcommand{\CharTok}[1]{\textcolor[rgb]{0.13,0.47,0.30}{#1}}
\newcommand{\CommentTok}[1]{\textcolor[rgb]{0.37,0.37,0.37}{#1}}
\newcommand{\CommentVarTok}[1]{\textcolor[rgb]{0.37,0.37,0.37}{\textit{#1}}}
\newcommand{\ConstantTok}[1]{\textcolor[rgb]{0.56,0.35,0.01}{#1}}
\newcommand{\ControlFlowTok}[1]{\textcolor[rgb]{0.00,0.23,0.31}{\textbf{#1}}}
\newcommand{\DataTypeTok}[1]{\textcolor[rgb]{0.68,0.00,0.00}{#1}}
\newcommand{\DecValTok}[1]{\textcolor[rgb]{0.68,0.00,0.00}{#1}}
\newcommand{\DocumentationTok}[1]{\textcolor[rgb]{0.37,0.37,0.37}{\textit{#1}}}
\newcommand{\ErrorTok}[1]{\textcolor[rgb]{0.68,0.00,0.00}{#1}}
\newcommand{\ExtensionTok}[1]{\textcolor[rgb]{0.00,0.23,0.31}{#1}}
\newcommand{\FloatTok}[1]{\textcolor[rgb]{0.68,0.00,0.00}{#1}}
\newcommand{\FunctionTok}[1]{\textcolor[rgb]{0.28,0.35,0.67}{#1}}
\newcommand{\ImportTok}[1]{\textcolor[rgb]{0.00,0.46,0.62}{#1}}
\newcommand{\InformationTok}[1]{\textcolor[rgb]{0.37,0.37,0.37}{#1}}
\newcommand{\KeywordTok}[1]{\textcolor[rgb]{0.00,0.23,0.31}{\textbf{#1}}}
\newcommand{\NormalTok}[1]{\textcolor[rgb]{0.00,0.23,0.31}{#1}}
\newcommand{\OperatorTok}[1]{\textcolor[rgb]{0.37,0.37,0.37}{#1}}
\newcommand{\OtherTok}[1]{\textcolor[rgb]{0.00,0.23,0.31}{#1}}
\newcommand{\PreprocessorTok}[1]{\textcolor[rgb]{0.68,0.00,0.00}{#1}}
\newcommand{\RegionMarkerTok}[1]{\textcolor[rgb]{0.00,0.23,0.31}{#1}}
\newcommand{\SpecialCharTok}[1]{\textcolor[rgb]{0.37,0.37,0.37}{#1}}
\newcommand{\SpecialStringTok}[1]{\textcolor[rgb]{0.13,0.47,0.30}{#1}}
\newcommand{\StringTok}[1]{\textcolor[rgb]{0.13,0.47,0.30}{#1}}
\newcommand{\VariableTok}[1]{\textcolor[rgb]{0.07,0.07,0.07}{#1}}
\newcommand{\VerbatimStringTok}[1]{\textcolor[rgb]{0.13,0.47,0.30}{#1}}
\newcommand{\WarningTok}[1]{\textcolor[rgb]{0.37,0.37,0.37}{\textit{#1}}}

\usepackage{longtable,booktabs,array}
\newcounter{none} % for unnumbered tables
\usepackage{calc} % for calculating minipage widths
% Correct order of tables after \paragraph or \subparagraph
\usepackage{etoolbox}
\makeatletter
\patchcmd\longtable{\par}{\if@noskipsec\mbox{}\fi\par}{}{}
\makeatother
% Allow footnotes in longtable head/foot
\IfFileExists{footnotehyper.sty}{\usepackage{footnotehyper}}{\usepackage{footnote}}
\makesavenoteenv{longtable}
\usepackage{graphicx}
\makeatletter
\newsavebox\pandoc@box
\newcommand*\pandocbounded[1]{% scales image to fit in text height/width
  \sbox\pandoc@box{#1}%
  \Gscale@div\@tempa{\textheight}{\dimexpr\ht\pandoc@box+\dp\pandoc@box\relax}%
  \Gscale@div\@tempb{\linewidth}{\wd\pandoc@box}%
  \ifdim\@tempb\p@<\@tempa\p@\let\@tempa\@tempb\fi% select the smaller of both
  \ifdim\@tempa\p@<\p@\scalebox{\@tempa}{\usebox\pandoc@box}%
  \else\usebox{\pandoc@box}%
  \fi%
}
% Set default figure placement to htbp
\def\fps@figure{htbp}
\makeatother





\setlength{\emergencystretch}{3em} % prevent overfull lines

\providecommand{\tightlist}{%
  \setlength{\itemsep}{0pt}\setlength{\parskip}{0pt}}



 


\usepackage{tcolorbox}
\usepackage{amssymb}
\usepackage{yfonts}
\usepackage{bm}


\newtcolorbox{greybox}{
  colback=white,
  colframe=blue,
  coltext=black,
  boxsep=5pt,
  arc=4pt}
  
\newcommand{\sectionbreak}{\clearpage}

 
\newcommand{\ds}[4]{\sum_{{#1}=1}^{#3}\sum_{{#2}=1}^{#4}}
\newcommand{\us}[3]{\mathop{\sum\sum}_{1\leq{#2}<{#1}\leq{#3}}}

\newcommand{\ol}[1]{\overline{#1}}
\newcommand{\ul}[1]{\underline{#1}}

\newcommand{\amin}[1]{\mathop{\text{argmin}}_{#1}}
\newcommand{\amax}[1]{\mathop{\text{argmax}}_{#1}}

\newcommand{\ci}{\perp\!\!\!\perp}

\newcommand{\mc}[1]{\mathcal{#1}}
\newcommand{\mb}[1]{\mathbb{#1}}
\newcommand{\mf}[1]{\mathfrak{#1}}

\newcommand{\eps}{\epsilon}
\newcommand{\lbd}{\lambda}
\newcommand{\alp}{\alpha}
\newcommand{\df}{=:}
\newcommand{\am}[1]{\mathop{\text{argmin}}_{#1}}
\newcommand{\ls}[2]{\mathop{\sum\sum}_{#1}^{#2}}
\newcommand{\ijs}{\mathop{\sum\sum}_{1\leq i<j\leq n}}
\newcommand{\jis}{\mathop{\sum\sum}_{1\leq j<i\leq n}}
\newcommand{\sij}{\sum_{i=1}^n\sum_{j=1}^n}
	
\KOMAoption{captions}{tableheading}
\makeatletter
\@ifpackageloaded{caption}{}{\usepackage{caption}}
\AtBeginDocument{%
\ifdefined\contentsname
  \renewcommand*\contentsname{Table of contents}
\else
  \newcommand\contentsname{Table of contents}
\fi
\ifdefined\listfigurename
  \renewcommand*\listfigurename{List of Figures}
\else
  \newcommand\listfigurename{List of Figures}
\fi
\ifdefined\listtablename
  \renewcommand*\listtablename{List of Tables}
\else
  \newcommand\listtablename{List of Tables}
\fi
\ifdefined\figurename
  \renewcommand*\figurename{Figure}
\else
  \newcommand\figurename{Figure}
\fi
\ifdefined\tablename
  \renewcommand*\tablename{Table}
\else
  \newcommand\tablename{Table}
\fi
}
\@ifpackageloaded{float}{}{\usepackage{float}}
\floatstyle{ruled}
\@ifundefined{c@chapter}{\newfloat{codelisting}{h}{lop}}{\newfloat{codelisting}{h}{lop}[chapter]}
\floatname{codelisting}{Listing}
\newcommand*\listoflistings{\listof{codelisting}{List of Listings}}
\makeatother
\makeatletter
\makeatother
\makeatletter
\@ifpackageloaded{caption}{}{\usepackage{caption}}
\@ifpackageloaded{subcaption}{}{\usepackage{subcaption}}
\makeatother
\usepackage{bookmark}
\IfFileExists{xurl.sty}{\usepackage{xurl}}{} % add URL line breaks if available
\urlstyle{same}
\hypersetup{
  pdftitle={Voronoi Homogeneity Analysis of Categorical Data},
  pdfauthor={Jan de Leeuw},
  colorlinks=true,
  linkcolor={blue},
  filecolor={Maroon},
  citecolor={Blue},
  urlcolor={Blue},
  pdfcreator={LaTeX via pandoc}}


\title{Voronoi Homogeneity Analysis of Categorical Data}
\author{Jan de Leeuw}
\date{May 3, 2026}
\begin{document}
\maketitle
\begin{abstract}
TBD
\end{abstract}

\renewcommand*\contentsname{Table of contents}
{
\setcounter{tocdepth}{3}
\tableofcontents
}

\sectionbreak

\textbf{Note:} This is a working manuscript which will be
expanded/updated frequently. All suggestions for improvement are
welcome. All Rmd, tex, html, pdf, R, and C files are in the public
domain. Attribution will be appreciated, but is not required. The files
can be found at \url{https://github.com/deleeuw/voronoi}

\sectionbreak

\section{Introduction}\label{introduction}

\section{Homogeneity Analysis}\label{homogeneity-analysis}

\sectionbreak

\section{Voronoi Analysis}\label{voronoi-analysis}

\[
\sigma(X,Y,\Delta)=\sum_{i=1}^n\sum_{j=1}^m\|\delta _{ij}-d_{ij}(X,Y)\|^2
\]

If \(g_{i\ell}^j=1\) then
\(\{\delta_{ij}\}_\ell\leq\{\delta_{ij}\}_\nu\) for all
\(\nu=1,\cdots,k_j\).

\[
\delta_{ij}'J\delta_{ij}=1
\] \# Majorization

\[
A_{ij}:=\begin{bmatrix}e_ie_i'&-e_ie_j'\\-e_je_i'&e_je_j'\end{bmatrix}
\] \[
Z:=\begin{bmatrix}
X\\Y
\end{bmatrix}
\]

\[
d_{ij}(X,Y)=\text{tr}\ Z'A_{ij}Z
\] \[
V:=\sum_{i=1}^n\sum_{j=1}^m A_{ij}=\begin{bmatrix}\hfill mI&-ee'\\-ee'&\hfill nI\end{bmatrix}
\]

If we define \begin{align}
V_1&:=(n+m)^{-1}\begin{bmatrix}\frac{m}{n}E&-E\\-E&\frac{n}{m}E\end{bmatrix},\\
V_2&:=\begin{bmatrix}J_n&0\\0&0\end{bmatrix},\\
V_3&:=\begin{bmatrix}0&0\\0&J_m\end{bmatrix},
\end{align} then \[
V=(n+m)V_1+mV_2+nV_3
\] Matrices \(V_1,V_2,\) and \(V_3\) are symmetric projectors on
orthogonal subspaces. It follows that \(V^+\), the Moore-Penrose
inverse3 of \(V\), is \[
V^+=\frac{1}{n+m}V_1+\frac{1}{m}V_2+\frac{1}{n}V_3
\] \[
B(X,Y):=\begin{bmatrix}B_{11}(X,Y)&B_{12}(X,Y)\\B_{21}(X,Y)&B_{22}(X,Y)\end{bmatrix}
\] Now \(V_1B(X,Y)=0\) and thus
\[V^+B(X,Y)=\frac{1}{m}V_2B(X,Y)+\frac{1}{n}V_3B(X,Y)\]

\section{The Cone}\label{the-cone}

Consider the cone of all vectors \(x\) in \(\mathbb{R}^n\) for which
\(x_i\leq x_1\) for all \(i\neq 1\) and \(e'x=0\). The extreme rays of
the cone will be the \(n-1\) vectors of the form
\(\alpha e_i+\beta(e-e_i)\) with \(i\neq 1\). Thus element \(i\) is
equal to \(\alpha\) and all other elements (including \(x_1\)) are equal
to beta. Thus we must have \(\alpha>\beta\). The sum of the elements is
\(\alpha+(n-1)\beta\), which must be zero. Thus \(\alpha=-(n-1)\beta\),
which implies \(\beta<0\) and \(\alpha>0\). The extreme rays are
positive multiples of the vectors \(e_i-n^{-1}e\) with \(i\neq 1\) and
the cone is the set of all weighted sums with non-negative weights of
these \(n-1\) vectors. Here is a small example.

\begin{Shaded}
\begin{Highlighting}[]
\NormalTok{a }\OtherTok{\textless{}{-}} \SpecialCharTok{{-}}\FunctionTok{cbind}\NormalTok{(}\DecValTok{1}\NormalTok{, }\SpecialCharTok{{-}}\FunctionTok{diag}\NormalTok{(}\DecValTok{4}\NormalTok{))}
\NormalTok{b }\OtherTok{\textless{}{-}}\NormalTok{ (}\FunctionTok{diag}\NormalTok{(}\DecValTok{5}\NormalTok{)}\SpecialCharTok{{-}}\NormalTok{(}\DecValTok{1} \SpecialCharTok{/} \DecValTok{5}\NormalTok{))[, }\SpecialCharTok{{-}}\DecValTok{1}\NormalTok{]}
\FunctionTok{print}\NormalTok{(a)}
\end{Highlighting}
\end{Shaded}

\begin{verbatim}
     [,1] [,2] [,3] [,4] [,5]
[1,]   -1    1    0    0    0
[2,]   -1    0    1    0    0
[3,]   -1    0    0    1    0
[4,]   -1    0    0    0    1
\end{verbatim}

\begin{Shaded}
\begin{Highlighting}[]
\FunctionTok{print}\NormalTok{(b)}
\end{Highlighting}
\end{Shaded}

\begin{verbatim}
     [,1] [,2] [,3] [,4]
[1,] -0.2 -0.2 -0.2 -0.2
[2,]  0.8 -0.2 -0.2 -0.2
[3,] -0.2  0.8 -0.2 -0.2
[4,] -0.2 -0.2  0.8 -0.2
[5,] -0.2 -0.2 -0.2  0.8
\end{verbatim}

\begin{Shaded}
\begin{Highlighting}[]
\FunctionTok{print}\NormalTok{(a }\SpecialCharTok{\%*\%}\NormalTok{ b)}
\end{Highlighting}
\end{Shaded}

\begin{verbatim}
     [,1] [,2] [,3] [,4]
[1,]    1    0    0    0
[2,]    0    1    0    0
[3,]    0    0    1    0
[4,]    0    0    0    1
\end{verbatim}

For later reference the sum of squares of the vectors \(e_i-n^{-1}e\) is
\((n-1)/n\), so the vectors \(\sqrt{n/(n-1)}(e_i-n^{-1}e)\) are the
vectors of length 1 on the extreme rays. The average of the extreme rays
is

\section{Normalized Monotone
Regression}\label{normalized-monotone-regression}

Suppose \(y\) is a vector with \(n\) elements. The problem we study is
to minimixe the sum of squares \(\sigma(x):=\|x-y\|^2\) over \(x\in K\)
and \(x'Jx=1\). Here

\begin{itemize}
\tightlist
\item
  \(K\) is the cone of vectors with
  \(\smash{x_1\leq x_2\leq\cdots\leq x_n}\),
\item
  \(J\) is the centering matrix \(\smash{J:=I-n^{-1}ee'}\), where
\item
  \(e\) is the vector with all \(n\) elements equal to one.
\end{itemize}

Change variables to \(\tilde x=Jx\). Then \(x'Jx=\tilde x'\tilde x\) and
\(x\in K\) if and only if \(\tilde x\in K\). Also \(\tilde x = Jx\) if
and only if \(x=\tilde x+\beta e\). So the transformed problem is to
minimize \(\|\tilde x+\beta e-y\|^2\) over \(\beta\) and
\(\tilde x\in K\), with \(e'\tilde x=0\) and \(\tilde x'\tilde x=1\).
The minimizer for \(\beta\) is \(\overline{y}:=n^{-1}e'y\), and thus the
problem becomes minimizing \(\|\tilde x- Jy\|^2\) over
\(\tilde x\in K\), \(\tilde x'\tilde x=1\), and \(e'\tilde x=0\).

Forget about the constraint \(e'\tilde x=0\). If \(M\) is the monotone
regression operator it follows that the solution for \(\tilde x\) is
\(M(Jy)/\|M(Jy)\|\), which automatically satisfies \(e'\tilde x=0\).
Thus the solution of our original problem is \[
\hat x=\frac{M(y-\overline{y})}{\|M(y-\overline{y})\|}+\overline{y}.
\] \(M(y-\overline{y})=M(y)-\overline{y}\) and thus
\(\|M(y-\overline y)\|^2=\|M(y)\|^2+n\overline{y}^2-2\overline{y}e'M(y)\)
nor \(e'M(y)=e'y=n\overline{y}\). Thus
\(\|M(y-\overline y)\|^2=\|M(y)\|^2-n\overline{M(y)}^2\) \[
\hat x=\frac{M(y)-\overline{y}e}{\|M(y-\overline{y})\|}+\overline{y}e=
\frac{1}{\|M(y-\overline{y})\|}M(y)+\left\{1-\frac{1}{\|M(y-\overline{y})\|}\right\}\overline{y}e
\]

\sectionbreak

\section{References}\label{references}




\end{document}
